\subsection{Structural Stability of Projected Descriptions}
  \label{subsec:noninjective-projection-entanglement}

  Entanglement, measurement, and decay correspond to distinct stability
  regimes of projected descriptions.
  Entanglement corresponds to the regime in which a non-factorizable
  projected description remains structurally stable.
  Measurement and decoherence correspond to selective stabilization:
  interaction with an environment amplifies certain relational features
  while rendering alternative descriptions inadmissible.
  Particle decay represents the regime in which the non-factorizable
  description becomes unstable, and stability is recovered through
  factorization into several localized configurations.

  \begin{figure}[t]
    \centering
    \begin{tikzpicture}[
      node distance=10mm and 18mm,
      box/.style={draw, rounded corners, align=center,
      inner sep=6pt},
      arr/.style={->, thick},
      lab/.style={font=\small, align=center}
    ]
    \node[box] (chi)
    {$\chi$\\[-2pt]\footnotesize coherence space\\
    [-2pt]\footnotesize (admissibility constraints)};
    \node[box, below=of chi] (proj)
    {non-injective\\[-2pt]\footnotesize projection $\Pi$};
    \node[box, below=of proj] (chieff)
    {$\chi_{\mathrm{eff}}$\\[-2pt]\footnotesize constrained
    reconstruction};
    \node[box, below left=16mm and 24mm of chieff] (stable)
    {stable\\[-2pt]\footnotesize (non-factorizable)};
    \node[box, below right=16mm and 24mm of chieff] (unstable)
    {unstable\\[-2pt]\footnotesize (non-factorizable)};
    \node[box, below=of stable] (ent)
    {entanglement\\[-2pt]\footnotesize resolution event};
    \node[box, below=of unstable] (frag)
    {fragmentation\\[-2pt]\footnotesize (decay)};
    \node[box, below=of frag] (sum)
    {$\chi_{\mathrm{eff},1}\oplus
    \chi_{\mathrm{eff},2}\oplus\cdots$\\
    [-2pt]\footnotesize factorizable descriptions};
    \draw[arr] (chi) -- (proj);
    \draw[arr] (proj) -- (chieff);
    \draw[arr] (chieff) -- (stable);
    \draw[arr] (chieff) -- (unstable);
    \draw[arr] (stable) -- (ent);
    \draw[arr] (unstable) -- (frag);
    \draw[arr] (frag) -- (sum);
    \end{tikzpicture}
    \caption{Reprojection branching induced by the finite-resolution
    projection~$\Pi$.
      $\chi$ is a coherence space whose admissibility structure constrains
      which reconstructions can consistently emerge; it does not pre-select
      a unique substrate configuration.
      At each reprojection step the effective state $\chi_{\mathrm{eff}}$
      either remains jointly non-factorizable---entanglement, arising when
      two degrees of freedom share a resolution cell---or fragments into
      localized factorizable descriptions (particle decay).
      The branching is a property of the reprojection dynamics, not of a
      distinguished substrate configuration.}
    \label{fig:noninjective-branching-entanglement-decay}
  \end{figure}

  A technical analysis is provided in
  Appendix~\ref{subsec:non-injective-projection-and-structural-factorization}.
